\subsection{Trace and Determinant}

% **Role:** quantitative matrix behavior and computation.
%
% Trace, determinant, rank, spectral radius, PSD/PD matrices, idempotents, nilpotents.

Matrix classes package operator structure in coordinates. Trace, determinant, rank, positivity, idempotents, nilpotents, spectral radius, and structured products give computable invariants.

\begin{definition} \label{def:matrix-trace-and-determinant} \hlt{Trace and Determinant of a Matrix}\\
For $A=[A_{j,k}]\in\operatorname{Mat}(n,n,\mathbb{F})$, the \hlt{trace} of $A$ is $\operatorname{tr}(A)\defeq A_{1,1}+\cdots+A_{n,n}$.\\
The \hlt{determinant} of $A$ is $\det A \defeq \sum\limits_{\sigma\in\mathfrak{S}_n} \sign(\sigma)A_{\sigma(1),1}\cdots A_{\sigma(n),n}$.
\end{definition}

\begin{definition} \label{def:gram-determinant} \hlt{Gram Determinant}\\
For vectors $\mathbf{u}_1,\ldots,\mathbf{u}_n$ in an inner-product space, define their \hlt{Gram determinant} by $\operatorname{GramDet}(\mathbf{u}_1,\ldots,\mathbf{u}_n) \defeq \det \operatorname{Gram}(\mathbf{u}_1,\ldots,\mathbf{u}_n)$.
\end{definition}

\begin{example} \label{ex:frobenius-trace-riesz-vector} \hlt{Trace Functional and the Frobenius Inner Product}\\
On $\operatorname{Mat}(n,n,\mathbb{F})$ with the Frobenius inner product from Example \ref{ex:function-and-matrix-inner-products}, the Riesz vector of the trace functional $A\mapsto\operatorname{tr}(A)$ is $I_n$, since $\innerprod[F]{A,I_n}=\sum\limits_{j=1}^{n}a_{jj}=\operatorname{tr}(A)$.
\end{example}

\begin{proposition} \label{prop:trace-arithmetic} \hlt{Trace Arithmetic}\\
Let $A,B\in\operatorname{Mat}(n,n,\mathbb{F})$ and $r\in\mathbb{F}$. Then
\begin{enumerate}[label=(\roman*)]
\setlength{\itemsep}{0pt}
\item $\operatorname{tr}(rA)=r\operatorname{tr}(A)$;
\item $\operatorname{tr}(A+B)=\operatorname{tr}(A)+\operatorname{tr}(B)$.
\end{enumerate}
\begin{proof}
We prove the two identities separately.
\begin{enumerate}[label=(\roman*)]
\setlength{\itemsep}{0pt}
\item Summing the diagonal entries of $rA$ gives $\operatorname{tr}(rA)=\sum\limits_{j=1}^{n}ra_{jj}=r\sum\limits_{j=1}^{n}a_{jj}=r\operatorname{tr}(A)$.
\item Summing the diagonal entries of $A+B$ gives $\operatorname{tr}(A+B)=\sum\limits_{j=1}^{n}(a_{jj}+b_{jj}) =\sum\limits_{j=1}^{n}a_{jj}+\sum_{j=1}^{n}b_{jj} =\operatorname{tr}(A)+\operatorname{tr}(B)$.
\end{enumerate}
\hlt{Structural Commentary}: Trace is linear because it is the sum of diagonal entries.\\
\hlt{Key Insight}: Scalar multiplication and addition pass through the diagonal sum.
\end{proof}
\end{proposition}

\begin{proposition} \label{prop:trace-cyclicity-two-factors} \hlt{Trace Is Cyclic for Two Factors}\\
If $A\in\operatorname{Mat}(m,n,\mathbb{F})$ and $B\in\operatorname{Mat}(n,m,\mathbb{F})$, then $\operatorname{tr}(AB)=\operatorname{tr}(BA)$.
\begin{proof}
Expand both traces by entries: $\operatorname{tr}(AB) =\sum\limits_{j=1}^{m}\sum\limits_{k=1}^{n}A_{j,k}B_{k,j} =\sum\limits_{k=1}^{n}\sum\limits_{j=1}^{m}B_{k,j}A_{j,k} =\operatorname{tr}(BA)$.\\
\hlt{Structural Commentary}: Cyclicity is just reindexing the same double sum.\\
\hlt{Key Insight}: The two double sums contain exactly the same scalar terms.
\end{proof}
\end{proposition}

\begin{proposition} \label{prop:trace-cyclic-products} \hlt{Cyclic Products}\\
Let $m\geq 2$, $A_j\in\operatorname{Mat}(n_j,n_{j+1},\mathbb{F})$ for $j=1,\ldots,m$, $n_{m+1}=n_1$, $P_{p,q}\defeq A_pA_{p+1}\cdots A_q$.\\
Then $\operatorname{tr}(P_{1,m})=\operatorname{tr}(A_mP_{1,m-1})$. Consequently, for $A,B,C\in\operatorname{Mat}(n,n,\mathbb{F})$, $\operatorname{tr}(ABC)=\operatorname{tr}(BCA)=\operatorname{tr}(CAB)$.
\begin{proof}
Since $P_{1,m-1}\in\operatorname{Mat}(n_1,n_m,\mathbb{F})$ and $A_m\in\operatorname{Mat}(n_m,n_1,\mathbb{F})$, by Proposition \ref{prop:trace-cyclicity-two-factors}, $\operatorname{tr}(P_{1,m}) =\operatorname{tr}(P_{1,m-1}A_m) =\operatorname{tr}(A_mP_{1,m-1})$.\\
Applying this cyclic move twice to $ABC$ gives the three-factor identity.\\
\hlt{Structural Commentary}: Trace preserves cyclic order, not arbitrary order.\\
\hlt{Key Insight}: One cyclic move is $\operatorname{tr}(XY)=\operatorname{tr}(YX)$ with $X=P_{1,m-1}$ and $Y=A_m$.
\end{proof}
\end{proposition}

\begin{example} \label{ex:trace-cyclicity-not-arbitrary-permutation} \hlt{Cyclicity Is Not Arbitrary Commutativity}\\
Let $A=E_{1,2}$, $B=E_{2,1}$, and $C=E_{1,1}$ in $\operatorname{Mat}(2,2,\mathbb{F})$. Then $\operatorname{tr}(ABC)=1$, while $\operatorname{tr}(ACB)=0$.\\
Thus trace permits cyclic rotations of factors, but not arbitrary rearrangements.
\end{example}

\begin{theorem} \label{thm:determinant-basic-laws} \hlt{Basic Determinant Laws}\\
Let $A,B\in\operatorname{Mat}(n,n,\mathbb{F})$.
\begin{enumerate}[label=(\roman*)]
\setlength{\itemsep}{0pt}
\item $\det I_n=1$.
\item $\det(A)$ is linear as a function of each column of $A$ when the other columns are fixed.
\item Swapping two columns negates the determinant; equal columns force $\det A=0$.
\item If $A$ is triangular, then $\det A=A_{1,1}\cdots A_{n,n}$.
\item $\det(AB)=\det(A)\det(B)$.
\item $A$ is invertible if and only if $\det A\neq0$. In that case $\det(A^{-1})=(\det A)^{-1}$.
\item $\det(A^t)=\det(A)$.
\end{enumerate}
\begin{proof}
\begin{enumerate}[label=(\roman*)]
\setlength{\itemsep}{0pt}
\item $\det I_n=\sum\limits_{\sigma\in\mathfrak{S}_n}\sign(\sigma)(I_n)_{\sigma(1),1}\cdots (I_n)_{\sigma(n),n}=1$ as product is nonzero only when $\sigma(k)=k \ \forall k$.
\item For fixed $k$, $\det(A_{\cdot,1},\ldots,\mathbf{x},\ldots,A_{\cdot,n}) =\sum\limits_{\sigma\in\mathfrak{S}_n} \sign(\sigma)x_{\sigma(k)}\prod\limits_{\ell\neq k}A_{\sigma(\ell),\ell}$, which is linear in $\mathbf{x}$.
\item If $B_{\cdot,k}=A_{\cdot,\tau(k)}$ for transposition $\tau$, then $\det B=\sum\limits_{\rho\in\mathfrak{S}_n}\sign(\rho\tau)\prod\limits_{k=1}^n A_{\rho(k),k} =-\det A$.\\
If the swapped columns are equal, then $B=A$, so $\det A=-\det A$ and $\det A=0$.
\item If $A$ is upper triangular and $\prod_{k=1}^n A_{\sigma(k),k}\neq0$, then $\sigma(k)\leq k$ for all $k$, hence $\sigma(k)=k$ for all $k$.\\
Thus $\det A=A_{1,1}\cdots A_{n,n}$. If $A$ is lower triangular, nonzero products require $\sigma(k)\geq k$ for all $k$, again forcing $\sigma(k)=k$ for all $k$, so the same diagonal product formula holds.
\item If $F$ is alternating and multilinear in columns, then expanding each column in the standard basis gives\\
$F(A_{\cdot,1},\ldots,A_{\cdot,n}) =\sum\limits_{i_1,\ldots,i_n}A_{i_1,1}\cdots A_{i_n,n}F(\mathbf{e}_{i_1},\ldots,\mathbf{e}_{i_n}) =\sum\limits_{\sigma\in\mathfrak{S}_n}\sign(\sigma)F(I_n)A_{\sigma(1),1}\cdots A_{\sigma(n),n}$.\\
The terms with repeated $i_j$ vanish by alternation, and the surviving ordered bases differ from $(\mathbf{e}_1,\ldots,\mathbf{e}_n)$ by sign of permutation. For $H(B)=\det(AB)$, map $H$ is alternating multilinear and $H(I_n)=\det A$.\\
Hence $\det(AB)=H(B)=H(I_n)\det B=\det(A)\det(B)$.
\item If $A$ is invertible, then $1=\det I_n=\det(AA^{-1})=\det(A)\det(A^{-1})$.\\
If $A_{\cdot,k}=\sum\limits_{\ell\neq k}c_\ell A_{\cdot,\ell}$, then $\det A=\sum\limits_{\ell\neq k}c_\ell\det(A_{\cdot,1},\ldots,A_{\cdot,\ell},\ldots,A_{\cdot,\ell},\ldots,A_{\cdot,n})=0$.\\
Thus $\det A\neq0$ iff the columns of $A$ are independent, iff $A$ is invertible.
\item Using the permutation definition and replacing $\sigma$ by $\sigma^{-1}$,\\
$\det(A^t) =\sum\limits_{\sigma\in\mathfrak{S}_n}\sign(\sigma)A_{1,\sigma(1)}\cdots A_{n,\sigma(n)} =\sum\limits_{\tau\in\mathfrak{S}_n}\sign(\tau)A_{\tau(1),1}\cdots A_{\tau(n),n} =\det(A)$; $\sign(\tau^{-1})=\sign(\tau)$.
\end{enumerate}
\hlt{Structural Commentary}: The determinant is the normalized alternating multilinear measure of columns.\\
\hlt{Key Insight}: Multilinearity expands columns; alternation kills repeated directions.
\end{proof}
\end{theorem}

\begin{figure}[H]
  \centering
  \begin{tikzpicture}[
  >=Stealth,
  every node/.style={font=\small},
  axis/.style={->, gray!70},
  vec/.style={->, line width=1pt},
  area/.style={fill=blue!12, draw=blue!45}
]
  \draw[axis] (-0.17,0) -- (2.75,0);
  \draw[axis] (0,-0.17) -- (0,1.90);

  \coordinate (O) at (0,0);
  \coordinate (u) at (2.06,0.39);
  \coordinate (v) at (0.69,1.50);
  \coordinate (uv) at ($(u)+(v)$);

  \draw[area] (O) -- (u) -- (uv) -- (v) -- cycle;
  \draw[vec, blue!75!black] (O) -- (u) node[midway, below=3pt] {$\mathbf{a}_1$};
  \draw[vec, red!75!black] (O) -- (v) node[midway, left=2pt] {$\mathbf{a}_2$};
  \draw[gray!55] (u) -- (uv) -- (v);
\end{tikzpicture}

  \caption{Determinant is signed column-volume; swaps reverse orientation and dependence gives zero volume.}
  \label{fig:determinant-column-geometry}
\end{figure}

\begin{method} \label{meth:determinant-by-structure} \hlt{Determinant by Structure}\\
To compute or simplify a determinant:
\begin{enumerate}[label=(\roman*)]
\setlength{\itemsep}{0pt}
\item first look for triangular or block triangular form;
\item use column linearity to split complicated columns only when it reduces structure;
\item use column swaps to move useful pivots into position, tracking a factor of $-1$ for each swap;
\item detect zero determinants from dependent columns or repeated columns;
\item for products or changes of basis, use multiplicativity instead of expanding permutations.
\end{enumerate}
These steps use Theorem \ref{thm:determinant-basic-laws}.
\end{method}

\begin{figure}[H]
  \centering
  \begin{tikzpicture}[
  >=Stealth,
  every node/.style={font=\small},
  axis/.style={->, gray!70},
  vec/.style={->, line width=1pt},
  areaA/.style={fill=blue!12, draw=blue!50},
  areaB/.style={fill=red!12, draw=red!55},
  arrow/.style={->, line width=0.85pt, blue!70!black}
]
  \begin{scope}
    \draw[axis] (-0.2,0) -- (2.4,0);
    \draw[axis] (0,-0.2) -- (0,2.0);
    \draw[areaA] (0,0) rectangle (1.35,1.05);
    \node at (0.68,0.52) {$\Omega$};
  \end{scope}

  \draw[arrow] (2.7,0.75) -- (4.1,0.75) node[midway, above=4pt] {$T$};

  \begin{scope}[shift={(4.5,0)}]
    \draw[axis] (-0.2,0) -- (3.0,0);
    \draw[axis] (0,-0.2) -- (0,2.2);
    \coordinate (O) at (0,0);
    \coordinate (a) at (2.25,0.35);
    \coordinate (b) at (0.65,1.55);
    \coordinate (ab) at ($(a)+(b)$);
    \draw[areaB] (O) -- (a) -- (ab) -- (b) -- cycle;
    \draw[vec, blue!75!black] (O) -- (a);
    \draw[vec, red!75!black] (O) -- (b);
    \node at (1.45,0.95) {$T(\Omega)$};
  \end{scope}
\end{tikzpicture}

  \caption{A linear map scales oriented volume by its determinant; singular maps collapse volume to zero.}
  \label{fig:matrix-determinant-volume-scaling}
\end{figure}

\begin{lemma} \label{lem:determinant-identity-direct-sum} \hlt{Identity Blocks Do Not Change Determinant}\\
Let $A\in\operatorname{Mat}(m,m,\mathbb{F})$ and let $n\geq0$. Then $\det(I_n\oplus A)=\det(A)=\det(A\oplus I_n)$.
\begin{proof}
In the permutation formula for $I_n\oplus A$, a nonzero term must satisfy $\sigma(k)=k$ for $1\leq k\leq n$. Remaining choices are exactly permutations of $\{n+1,\ldots,n+m\}$, after relabeling, gives $\det(A)$. Hence $\det(I_n\oplus A)=\det(A)$.\\
The same argument applied to the last $n$ coordinates gives $\det(A\oplus I_n)=\det(A)$.\\
The case $n=0$ is the empty identity block convention.\\
\hlt{Structural Commentary}: Identity blocks add directions without changing volume.\\
\hlt{Key Insight}: In the determinant expansion, the identity block forces the matching coordinates to be fixed.
\end{proof}
\end{lemma}

\begin{theorem} \label{thm:block-upper-triangular-determinant} \hlt{Determinant of an Upper Block-Triangular Matrix}\\
Let $A\in\operatorname{Mat}(r,r,\mathbb{F})$, $B\in\operatorname{Mat}(r,s,\mathbb{F})$, and $D\in\operatorname{Mat}(s,s,\mathbb{F})$. Then $\det
\begin{bmatrix}
A & B\\
0 & D
\end{bmatrix}
=\det(A)\det(D)$.
\begin{proof}
Factor $\begin{bmatrix}
A & B\\
0 & D
\end{bmatrix}
=
\begin{bmatrix}
I_r & 0\\
0 & D
\end{bmatrix}
\begin{bmatrix}
I_r & B\\
0 & I_s
\end{bmatrix}
\begin{bmatrix}
A & 0\\
0 & I_s
\end{bmatrix}$.\\
The first and last factors have determinants $\det(D)$ and $\det(A)$ by Lemma \ref{lem:determinant-identity-direct-sum}.\\
The middle factor is upper triangular with every diagonal entry equal to $1$, so its determinant is $1$.\\
Multiplicativity of determinant gives the formula.\\
\hlt{Structural Commentary}: Block upper triangular matrices separate determinant contribution into diagonal blocks. The off-diagonal block does not affect the determinant.\\
\hlt{Key Insight}: Factor the matrix into two block diagonal pieces and one unit upper triangular piece.
\end{proof}
\end{theorem}

\begin{proposition} \label{prop:schur-complement-determinant-formula} \hlt{Schur Complement Determinant Formula}\\
Let $M=
\begin{bmatrix}
A & B\\
C & D
\end{bmatrix}$, $A\in\operatorname{Mat}(r,r,\mathbb{F})$, $D\in\operatorname{Mat}(s,s,\mathbb{F})$.\\
If $A$ is invertible, then the \hlt{Schur complement of $A$ in $M$} is $M/A\defeq D-CA^{-1}B$, and $\det(M)=\det(A)\det(M/A)$.
\begin{proof}
Block Gaussian elimination gives $\begin{bmatrix}
I_r & 0\\
-CA^{-1} & I_s
\end{bmatrix}
\begin{bmatrix}
A & B\\
C & D
\end{bmatrix}
=
\begin{bmatrix}
A & B\\
0 & D-CA^{-1}B
\end{bmatrix}$.\\
The elimination matrix is lower block-triangular with identity diagonal blocks, so its determinant is $1$. Take determinants, apply Theorem \ref{thm:block-upper-triangular-determinant} to right-hand block triangular matrix, $\det(M)=\det(A)\det(D-CA^{-1}B)$.\\
\hlt{Structural Commentary}: Schur complementation is block Gaussian elimination.\\
\hlt{Key Insight}: Eliminating $C$ splits the determinant into $\det(A)$ times the remaining block determinant.
\end{proof}
\end{proposition}

\begin{proposition} \label{prop:schur-complement-lower-right-determinant} \hlt{Schur Complement of the Lower-Right Block}\\
Let $M=
\begin{bmatrix}
A & B\\
C & D
\end{bmatrix}$, $A\in\operatorname{Mat}(r,r,\mathbb{F})$, $D\in\operatorname{Mat}(s,s,\mathbb{F})$.\\
If $D$ is invertible, then the \hlt{Schur complement of $D$ in $M$} is $M/D\defeq A-BD^{-1}C$, and $\det(M)=\det(D)\det(M/D)$.
\begin{proof}
Block Gaussian elimination gives $\begin{bmatrix}
I_r & -BD^{-1}\\
0 & I_s
\end{bmatrix}
\begin{bmatrix}
A & B\\
C & D
\end{bmatrix}
=
\begin{bmatrix}
A-BD^{-1}C & 0\\
C & D
\end{bmatrix}$.\\
The elimination matrix has determinant $1$. The transpose of the matrix on the right is upper block-triangular, and $\det(X^t)=\det(X)$ by Theorem \ref{thm:determinant-basic-laws}. By Theorem \ref{thm:block-upper-triangular-determinant}, $\det(M)=\det(A-BD^{-1}C)\det(D)$.\\
\hlt{Structural Commentary}: This is the lower-right version of Schur complementation. Eliminating the upper-right block isolates the determinant contribution of $D$.\\
\hlt{Key Insight}: Eliminate by blocks, then read the determinant from the diagonal blocks.
\end{proof}
\end{proposition}

\begin{proposition} \label{prop:bordered-determinant-formulas} \hlt{Bordered Determinant Formulas}\\
Let $A\in\operatorname{Mat}(n,n,\mathbb{F})$, let $\mathbf{x},\mathbf{y}\in\mathbb{F}^n$, and let $c\in\mathbb{F}$.
\begin{enumerate}[label=(\roman*)]
\setlength{\itemsep}{0pt}
\item If $c\neq0$, then $\det
\begin{bmatrix}
c & \mathbf{x}^t\\
\mathbf{y} & A
\end{bmatrix}
=
c\det\left(A-\frac{1}{c}\mathbf{y}\mathbf{x}^t\right)$.
\item If $A$ is invertible, then $\det
\begin{bmatrix}
A & \mathbf{x}\\
\mathbf{y}^t & c
\end{bmatrix}
=
\left(c-\mathbf{y}^tA^{-1}\mathbf{x}\right)\det(A)$
\end{enumerate}
\begin{proof}
\begin{enumerate}[label=(\roman*)]
\setlength{\itemsep}{0pt}
\item This is Proposition \ref{prop:schur-complement-determinant-formula} with the scalar block $c$ in the upper-left corner.
\item This is Proposition \ref{prop:schur-complement-determinant-formula} with $A$ as the upper-left block and $D=c$.
\end{enumerate}
\hlt{Structural Commentary}: Bordered determinants are Schur complements with scalar borders.\\
\hlt{Key Insight}: Treat the border as a $1\times1$ block and apply the Schur complement formula.
\end{proof}
\end{proposition}

\begin{theorem} \label{thm:sylvester-determinant-identity} \hlt{Sylvester Determinant Identity}\\
Let $A\in\operatorname{Mat}(n,m,\mathbb{F})$ and $B\in\operatorname{Mat}(m,n,\mathbb{F})$. Then $\det(I_n+AB)=\det(I_m+BA)$.
\begin{proof}
Let $M=
\begin{bmatrix}
I_n & -A\\
B & I_m
\end{bmatrix}$.\\
By Proposition \ref{prop:schur-complement-lower-right-determinant} with lower-right block $I_m$, $\det(M)=\det(I_n+AB)$. By Proposition \ref{prop:schur-complement-determinant-formula} with upper-left block $I_n$, $\det(M)=\det(I_m+BA)$. Hence the two determinants are equal.\\
\hlt{Structural Commentary}: Two Schur complements of one block matrix give Sylvester's identity.\\
\hlt{Key Insight}: Compare the same determinant after eliminating different identity blocks.
\end{proof}
\end{theorem}

\begin{theorem} \label{thm:ab-ba-nonzero-eigenvalues} \hlt{Nonzero Eigenvalues of $AB$ and $BA$}\\
Let $A\in\operatorname{Mat}(m,n,\mathbb{C})$, $B\in\operatorname{Mat}(n,m,\mathbb{C})$, and suppose $n\geq m$. Then $p_{BA}(z)=z^{n-m}p_{AB}(z)$.\\
Consequently, for $\lambda\neq0$, $m_{AB}(\lambda)=m_{BA}(\lambda)$. If $m_{AB}(0)=k$, then $m_{BA}(0)=k+n-m$.
\begin{proof}
For $z\neq0$, by Theorem \ref{thm:sylvester-determinant-identity} applied to $-z^{-1}B$ and $A$,
\begin{align}
\det(zI_n-BA) =z^n\det(I_n-z^{-1}BA) =z^n\det(I_m-z^{-1}AB) =z^{n-m}\det(zI_m-AB).
\nonumber
\end{align}
Both sides are polynomials in $z$, so the identity holds for all $z$. The claimed multiplicities follow by factoring $p_{AB}$ and comparing with $p_{BA}=z^{n-m}p_{AB}$.\\
\hlt{Structural Commentary}: Sylvester's determinant identity moves the rectangular product across the determinant.\\
\hlt{Key Insight}: Passing from $AB$ to $BA$ only adds the necessary extra zero roots when the sizes differ.
\end{proof}
\end{theorem}

\begin{theorem} \label{thm:similarity-preserves-trace-and-determinant} \hlt{Trace and Determinant Are Similarity Invariants}\\
If $A,B\in\operatorname{Mat}(n,n,\mathbb{F})$ are similar, then $\operatorname{tr}(A)=\operatorname{tr}(B)$, $\det(A)=\det(B)$.
\begin{proof}
Write $B=PAP^{-1}$ with $P$ invertible.\\
Proposition \ref{prop:trace-cyclicity-two-factors} gives $\operatorname{tr}(B)=\operatorname{tr}\bigl(P(AP^{-1})\bigr) =\operatorname{tr}\bigl((AP^{-1})P\bigr) =\operatorname{tr}(A)$.\\
Theorem \ref{thm:determinant-basic-laws} gives $\det(B)=\det(P)\det(A)\det(P^{-1})=\det(A)$.\\
\hlt{Structural Commentary}: Trace and determinant ignore the coordinate system used to write an operator.\\
\hlt{Key Insight}: Conjugating by $P$ inserts $P$ and $P^{-1}$, which cancel inside trace and determinant.
\end{proof}
\end{theorem}

\begin{perspective} \label{per:similarity-invariants-and-canonical-forms} \hlt{Similarity Invariants and Canonical Forms}\\
Similarity is equality up to change of basis. Trace, determinant, characteristic polynomials, and minimal polynomials are similarity invariants by Theorems \ref{thm:similarity-preserves-trace-and-determinant}, \ref{thm:similarity-preserves-characteristic-polynomial}, and \ref{thm:similarity-preserves-minimal-polynomial}, but they are not generally complete.\\
Canonical forms, such as the Jordan form of Theorem \ref{thm:complex-jordan-form}, choose one structured representative of each similarity class when possible.
\end{perspective}

\begin{principle} \label{princ:trace-determinant-coordinate-free} \hlt{Trace and Determinant Are Coordinate-Free}\\
Trace and determinant are first computed from a matrix, but their real structural meaning is operator-level: changing basis conjugates the matrix, and conjugation preserves both quantities.
\end{principle}

\begin{theorem} \label{thm:operator-matrix-trace-determinant-basis-independent} \hlt{Trace and Determinant of an Operator Are Basis-Independent}\\
Let $\mathcal{V}$ be finite-dimensional, let $T\in\LL(\mathcal{V})$, and let $\beta,\gamma$ be bases of $\mathcal{V}$.\\
Then $\operatorname{tr}\bigl([T]_{\beta\leftarrow\beta}\bigr) =\operatorname{tr}\bigl([T]_{\gamma\leftarrow\gamma}\bigr)$, $\det\bigl([T]_{\beta\leftarrow\beta}\bigr) =\det\bigl([T]_{\gamma\leftarrow\gamma}\bigr)$.
\begin{proof}
By Theorem \ref{thm:change-of-basis-formula}, the two matrices are similar: $[T]_{\gamma\leftarrow\gamma} =C_{\gamma\leftarrow\beta}[T]_{\beta\leftarrow\beta}C_{\beta\leftarrow\gamma}$.\\
Theorem \ref{thm:similarity-preserves-trace-and-determinant} now gives equality of traces and determinants.\\
\hlt{Structural Commentary}: Basis changes conjugate matrices; trace and determinant remain operator invariants.\\
\hlt{Key Insight}: Similar matrices have the same trace and determinant.
\end{proof}
\end{theorem}

\begin{figure}[H]
  \centering
  \begin{tikzpicture}[
  >=Stealth,
  every node/.style={font=\small},
  box/.style={draw=gray!55, rounded corners=3pt, minimum width=2.5cm, minimum height=0.84cm, align=center, fill=gray!4},
  map/.style={->, line width=0.8pt, blue!65!black}
]
  \node[box] (beta) at (0,0) {$[T]_{\beta\leftarrow\beta}=A$};
  \node[box] (gamma) at (4.2,0) {$[T]_{\gamma\leftarrow\gamma}=PAP^{-1}$};

  \draw[map] (beta) -- (gamma)
    node[midway, above] {$P=C_{\gamma\leftarrow\beta}$};
\end{tikzpicture}

  \caption{Basis changes conjugate matrices; trace and determinant are unchanged.}
  \label{fig:operator-trace-determinant-basis-invariance}
\end{figure}

\begin{definition} \label{def:operator-trace-and-determinant} \hlt{Trace and Determinant of an Operator}\\
Let $\mathcal{V}$ be finite-dim, let $T\in\LL(\mathcal{V})$. Choose any basis $\beta$ of $\mathcal{V}$. Define $\operatorname{tr}T\defeq\operatorname{tr}\bigl([T]_{\beta\leftarrow\beta}\bigr)$, $\det T\defeq\det\bigl([T]_{\beta\leftarrow\beta}\bigr)$.\\
Theorem \ref{thm:operator-matrix-trace-determinant-basis-independent} shows these definitions do not depend on chosen basis.
\end{definition}

\begin{proposition} \label{prop:isometry-determinant-absolute-one} \hlt{Unitary Matrices and Isometries Have Determinant of Modulus One}\\
Let $U\in\operatorname{Mat}(n,n,\mathbb{F})$ be unitary. Then $\abs{\det U}=1$. More generally, if $S\in\LL(\mathcal{V})$ is an isometry on a finite-dimensional inner-product space over $\mathbb{F}\in\{\mathbb{R},\mathbb{C}\}$, then $\abs{\det S}=1$.
\begin{proof}
The permutation definition gives $\det(U^*)=\overline{\det U}$. By Theorem \ref{thm:determinant-basic-laws}, taking determinants in $U^*U=I_n$ gives $1=\det(U^*)\det(U)=\overline{\det U}\det U=\abs{\det U}^2$. Hence $\abs{\det U}=1$.\\
For an isometry $S$, choose an orthonormal basis and let $U$ be the matrix of $S$ in that basis. Then $U$ is unitary by Theorem \ref{thm:finite-dimensional-isometry-iff-adjoint-product-identity}, and $\det S=\det U$.\\
\hlt{Structural Commentary}: In orthonormal coordinates, isometries are unitary or orthogonal, so the determinant calculation is the matrix one.\\
\hlt{Key Insight}: Determinants turn $U^*U=I$ into $\abs{\det U}^2=1$.
\end{proof}
\end{proposition}

\begin{proposition} \label{prop:operator-trace-determinant-characteristic-coefficients} \hlt{Trace and Determinant from $q_T$}\\
Let $\mathcal{V}\neq\{\mathbf{0}\}$ be finite-dim over $\mathbb{F}\in\{\mathbb{R},\mathbb{C}\}$, let $T\in\LL(\mathcal{V})$, $q_T(z)=z^n+c_{n-1}z^{n-1}+\cdots+c_1z+c_0$, $n=\dim\mathcal{V}$.\\
Then $\operatorname{tr}T=-c_{n-1}$ and $\det T=(-1)^nc_0$. If $\mathbb{F}=\mathbb{C}$ and $\lambda_1,\ldots,\lambda_n$ list the eigenvalues of $T$ with algebraic multiplicity, then $\operatorname{tr}T=\lambda_1+\cdots+\lambda_n$, $\det T=\lambda_1\cdots\lambda_n$.
\begin{proof}
Choose a basis $\beta$ and set $A=[T]_{\beta\leftarrow\beta}$.\\
Then $q_T=p_A$, and by Definition \ref{def:operator-trace-and-determinant}, $\operatorname{tr}T=\operatorname{tr}A$ and $\det T=\det A$.\\
Theorem \ref{thm:matrix-characteristic-polynomial-basic-coefficients} gives $\operatorname{tr}T=-c_{n-1}$ and $\det T=(-1)^nc_0$.\\
If $\mathbb{F}=\mathbb{C}$, then comparing $q_T(z)=(z-\lambda_1)\cdots(z-\lambda_n)$ with the displayed form of $q_T$ gives the formulas.\\
\hlt{Structural Commentary}: Characteristic polynomial stores trace and determinant in first and last coefficients.\\
\hlt{Key Insight}: Compare coefficients of $\det(zI-T)$.
\end{proof}
\end{proposition}

\begin{proposition} \label{prop:trace-determinant-from-eigenvalues} \hlt{Trace and Determinant from Eigenvalues}\\
Let $\lambda_1,\ldots,\lambda_n$ list the eigenvalues of $A\in\operatorname{Mat}(n,n,\mathbb{C})$ according to $m_A$. Then $\operatorname{tr}A=\sum\limits_{j=1}^{n}\lambda_j$ and $\det A=\prod\limits_{j=1}^{n}\lambda_j$.
\begin{proof}
By Fundamental Theorem of Algebra, $p_A(z)=\prod\limits_{j=1}^{n}(z-\lambda_j) =z^n-\left(\sum\limits_{j=1}^{n}\lambda_j\right)z^{n-1}+\cdots+(-1)^n\prod\limits_{j=1}^{n}\lambda_j$.\\
Theorem \ref{thm:matrix-characteristic-polynomial-basic-coefficients} gives $p_A(z)=z^n-(\operatorname{tr}A)z^{n-1}+\cdots+(-1)^n\det A$. Uniqueness of coefficients gives the identities.\\
\hlt{Structural Commentary}: Coefficients of $p_A$ encode symmetric functions of the eigenvalues.\\
\hlt{Key Insight}: Trace is the first root sum and determinant is the full root product.
\end{proof}
\end{proposition}

\begin{proposition} \label{prop:transpose-conjugate-adjoint-characteristic-polynomial} \hlt{Transpose, Conjugate, and Adjoint Spectra}\\
Let $A\in\operatorname{Mat}(n,n,\mathbb{C})$. Then $p_{A^t}(z)=p_A(z)$, $p_{\overline{A}}(z)=\overline{p_A(\overline{z})}$, and $p_{A^*}(z)=\overline{p_A(\overline{z})}$.\\
Consequently, if $\lambda\in\operatorname{spec}(A)$, then $m_{A^t}(\lambda)=m_A(\lambda)$ and $m_{\overline{A}}(\overline{\lambda})=m_{A^*}(\overline{\lambda})=m_A(\lambda)$.
\begin{proof}
Use determinant identities:
\begin{align}
p_{A^t}(z)&=\det(zI_n-A^t)=\det((zI_n-A)^t)=p_A(z), \nonumber\\
p_{\overline{A}}(z)&=\det(zI_n-\overline{A}) =\det(\overline{\overline{z}I_n-A}) =\overline{\det(\overline{z}I_n-A)} =\overline{p_A(\overline{z})}.
\nonumber
\end{align}
Since $A^*=\overline{A}^{\,t}$, the first identity applied to $\overline{A}$ gives the formula for $A^*$.\\
Multiplicity is preserved because conjugating a polynomial sends a factor $(z-\lambda)^k$ to $(z-\overline{\lambda})^k$.\\
\hlt{Structural Commentary}: Transpose preserves $p_A$; conjugation conjugates its roots.\\
\hlt{Key Insight}: The adjoint combines conjugation with transposition, so its spectrum is conjugated.
\end{proof}
\end{proposition}
